\chapter{The Gaussian case and the LMMSE baseline}
\label{ch:gaussian}

The Gaussian chain plays two roles: it is the case where everything is computable in closed
form, hence the yardstick for validating the general machinery; and it is the model implicitly
assumed by the most common approximation, hence the baseline everything is measured against.

\section{Everything in closed form}

Take $\innov = \normal{0}{\ivar}$. Then $a$ is jointly Gaussian with covariance
$\Sigma_0 = \corr^{|i-j|}$ from \eqref{eq:ar1-cov}, and by \eqref{eq:forward-solution} so is
$x$, with
\begin{equation}
\Sigma_t \;=\; e^{-2t}\Sigma_0 + \Dt I .
\end{equation}
For jointly Gaussian variables the conditional mean is linear, giving \eqref{eq:gaussian-mean}:
$\pmean{a} = e^{-t}\Sigma_0\Sigma_t^{-1}x$ and $\score = -\Sigma_t^{-1}x$.

\begin{intuition}
The score is a fixed linear map. No inference is required at all --- one matrix solve. This is
exactly why Gaussian data is a poor probe of how models exploit structure: the answer is
linear regression, and every method that can do linear regression is optimal.
\end{intuition}

\subsection{The precision matrix is tridiagonal}

A useful structural fact. The clean precision $Q_0 = \Sigma_0^{-1}$ of an AR(1) chain is
\emph{tridiagonal}: writing the density as
$\Pz(a)\propto\exp[-\tfrac{1}{2\ivar}\sum_i (a_i-\corr a_{i-1})^2]$ and expanding, only terms
$a_i^2$ and $a_ia_{i-1}$ appear. So
\[
(Q_0)_{ii} = \frac{1+\corr^2}{\ivar}\ \text{(interior)},
\qquad
(Q_0)_{i,i\pm1} = -\frac{\corr}{\ivar},
\qquad
(Q_0)_{ij}=0 \ \text{for } |i-j|>1 .
\]
The posterior precision adds the likelihood's contribution, $J = (e^{-2t}/\Dt)I + Q_0$, and
stays tridiagonal at every $t$.

\begin{intuition}
Tridiagonal precision is the algebraic image of the Markov property: zero in the precision
matrix means conditional independence given everything else. The \emph{covariance} $\Sigma_t$
is dense --- every site is correlated with every other --- while the precision is banded. Being
Markov is a statement about the precision, not the covariance.
\end{intuition}

\section{The moment closure}
\label{sec:closure-chapter}

For a non-Gaussian innovation the messages of Chapter~\ref{ch:bp} are genuinely non-Gaussian
functions. A natural approximation keeps only two numbers per message: replace each message by
the Gaussian with the same mean and variance. This is cheap, requires no grid, and is what one
would write first.

Under the transition $a_i = \corr a_{i-1}+\varepsilon_i$ with $\varepsilon$ of mean $0$ and
variance $\ivar$, a message with moments $(m,v)$ maps to
\begin{equation}
\label{eq:closure-both}
\text{forward:}\quad (m,v)\mapsto(\corr m,\ \corr^2 v + \ivar);
\qquad
\text{backward:}\quad (m,v)\mapsto(m/\corr,\ (v+\ivar)/\corr^2),
\end{equation}
the backward map being the forward one inverted, because that recursion integrates over the
\emph{output} of the transition rather than its input. It requires $\corr\neq0$.

\begin{proposition}
\label{prop:closure-lmmse}
For the chain \eqref{eq:ar1} with $\E[a]=0$ and Gaussian unary likelihoods, the moment-closed
recursion returns
\begin{equation}
\label{eq:lmmse-chapter}
\lmmse \;=\; e^{-t}\Sigma_0\left(e^{-2t}\Sigma_0+\Dt I\right)^{-1}x ,
\end{equation}
the LMMSE estimator of the covariance-matched Gaussian model, for \emph{every} innovation law
with those two moments.
\end{proposition}

\begin{proof}
Both maps in \eqref{eq:closure-both} depend on $\innov$ only through $(0,\ivar)$, and the
likelihood absorption is a Gaussian--Gaussian product, which is exact. So the recursion returns
the same function of $x$ for every innovation law with those moments --- in particular the same
as on the covariance-matched \emph{Gaussian} chain. On that chain every message is genuinely
Gaussian, so the projection is the identity, the recursion is exact, and it returns the true
posterior mean. For jointly Gaussian zero-mean $(a,x)$ the posterior mean is linear and
coincides with the LMMSE estimator of Proposition~\ref{prop:lmmse-general}.
\end{proof}

\begin{pitfall}
The zero-mean hypothesis is load-bearing. The general LMMSE estimator is
$\E[a]+\Cov(a,x)\Cov(x)^{-1}(x-\E[x])$, and dropping the mean terms is only valid when both
means vanish. With a prior mean of $1.3$, \eqref{eq:lmmse-chapter} departs from the true LMMSE
by $0.65$ in maximum absolute value.
\end{pitfall}

\begin{intuition}
This is a sharper statement than ``the Gaussian approximation is approximate''. The closure
does not merely lose some information about the innovation --- it loses \emph{all} of it beyond
the variance. Two chains from the matched family of \S\ref{sec:matched-covariance}, one
Gaussian and one heavily bimodal, are given identical treatment. The approximation is blind by
construction, not by degree.

A consequence worth stating: at the single-Gaussian level, ``error from representing messages
badly'' and ``error from assuming the wrong model'' are the same number. They are different
objects only for richer message families.
\end{intuition}

\begin{codebox}
\coderef{src/bp\_gaussian.py}{gaussian\_chain\_bp} implements the closure in information form,
carrying precision $\lambda$ and information $h$ rather than mean and variance. That
parametrisation matters: a message conveying no information is $\lambda=0$ exactly, whereas in
moment form it is a variance of $+\infty$.
\end{codebox}

\begin{pitfall}
An earlier version of this project implemented the closure by evaluating each Gaussian message
\emph{on the grid}, pushing it through the exact update, and re-fitting moments. At weakly
informative $t$ the outgoing message is nearly flat, and moment-matching a flat function
\emph{truncated} to $[-A,A]$ returns a variance of $(2A)^2/12$ rather than $\infty$ --- an
invented precision of order $3/A^2$. The error fed back through the recursion, driving messages
towards the boundary. Measured maximum absolute error in the posterior mean at $t=1.8$:
$9.1\times10^{-1}$, against $4.4\times10^{-16}$ for the information form. The lesson is that
the closure must be done analytically; representing an unnormalisable object on a finite grid
is not a neutral implementation detail.
\end{pitfall}

\section{Where the approximation costs something}

Since the closure is exactly LMMSE, the gap to exact inference is the value of the
higher-order structure. Measured on a Laplace chain at $\corr=0.85$, the relative deviation of
the closure's score from the exact one falls monotonically from $0.198$ at $t=0.08$ to
$9.9\times10^{-5}$ at $t=2.4$.

\begin{intuition}
The direction may be surprising: the approximation is \emph{worst at low noise}. The reason is
that at large $t$ the likelihood is weak and the posterior is dominated by the smoothing effect
of the noise, which Gaussianises everything --- so a Gaussian approximation of a nearly Gaussian
object is nearly exact. At small $t$ the posterior is sharply peaked and its shape is inherited
from the prior, whose non-Gaussianity is then fully visible.

Whether that matters for \emph{generation} is a separate question, because errors at small $t$
enter the reverse process where the denoiser is barely acting. That question is what the
experiments of Chapter~\ref{ch:experiments} address, and the answer is not the obvious one.
\end{intuition}

\begin{codebox}
\coderef{experiments/exp\_02\_laplace\_gaussian\_message\_error.py}{} produces the numbers
above; \coderef{experiments/exp\_03\_nongaussian\_innovation\_sweep.py}{} extends them across
the innovation families of \S\ref{sec:matched-covariance}.
\end{codebox}

\section*{Sources}
\addcontentsline{toc}{section}{Sources}
The linear-minimum-mean-square-error estimator and its orthogonality characterisation are
standard \citep{kay1993,lehmann1998}. The Gaussian message recursions coincide with Kalman
filtering and RTS smoothing \citep{kalman1960,rauch1965}.
